\section{The Target-Delay Anomaly}\label{sec:target-delay}

A natural engineering impulse, once the greedy chain of Model~4* is
in place, is to polish each robot's tour with a standard tour-length
optimizer such as 2-opt. Tours produced by greedy $p/d^2$ extension
are locally good at each step but globally sub-optimal as traveling
salesman tours; 2-opt will, on average, shorten them. Doing so
\emph{hurts} the finder competitive ratio.

At the default configuration ($n=30$, $R=3$, $E=14$, $L=10$; 1,000
paired trials; seed~42), 2-opt post-processing reduces mean finder
tour length from 12.567 to 12.024 units --- a mean reduction of
0.543 units per tour, exactly as expected from a tour-length
optimizer. Over the same 1,000 trials, mean finder competitive ratio
\emph{rises} from 3.093 under the greedy chain to 3.257 under the
2-opt chain, a 5.3 percent degradation. The effect is distributional
rather than universal: 225 of the 1,000 trials (22.5\%) are strictly
worse under 2-opt, 491 (49.1\%) are strictly better, and the
remaining 284 (28.4\%) are unchanged. More trials are helped than
hurt, but the trials that are hurt are hurt by enough to drag the
mean upward.

The mechanism is that 2-opt optimizes \emph{geometric} tour length,
a quantity that is blind to the prior probability $p_i$ attached to
each site. The greedy extension of Model~4* instead orders sites by
$p/d^2$: at every step it prefers a site that is both nearby and
high-probability. Because the target is drawn from the prior, it is
more likely to fall in a high-probability site than a low-probability
one, so the greedy chain tends to reach the target through a prefix
that accumulates high-probability-weighted distance efficiently ---
every leg of the prefix is either short or highly-rewarded by the
prior. 2-opt reorders the same set of sites to minimize total tour
length, but it does so without reference to the prior: sites are
swapped into earlier or later positions according to their geometry
alone. Two consequences follow. First, the rank at which the target
appears in the tour is reshuffled (the mean rank shifts only slightly,
from 3.220 to 3.203, but the distribution spreads: 2-opt has fewer
trials with the target at rank~3 and more at rank~5 and 7 ---
Figure~\ref{fig:target-delay}, left panel). Second, and more
importantly, even when the target occupies the same rank under both
orderings, the 2-opt prefix preceding it is composed of sites chosen
for geometric efficiency rather than $p/d^2$ value, which typically
lengthens the finder's distance-to-target. The scatter in
Figure~\ref{fig:target-delay} (right panel) makes the anomaly visible:
121 trials lie in the upper-right quadrant, where 2-opt shortened the
overall tour yet increased finder distance. The overall correlation
between tour-length reduction and FCR increase is weakly positive
(Pearson $r = 0.054$), consistent with the interpretation that 2-opt
and probability-weighted ordering are approximately orthogonal:
shortening the tour neither systematically helps nor systematically
hurts a finder-centric objective.

\begin{figure}[ht]
  \centering
  \includegraphics[width=0.95\textwidth]{fig_target_delay.png}
  \caption{The Target-Delay Anomaly on 1,000 paired trials at the
    default configuration. \textbf{Left:} histogram of the target's
    rank within the finder's tour for Model~4* under greedy chain
    extension (green) and the same tours post-processed with 2-opt
    (red). Mean ranks are nearly identical (3.220 vs.\ 3.203); 2-opt
    redistributes mass rather than shifting it wholesale. \textbf{Right:}
    per-trial scatter of tour-length reduction against FCR increase.
    Points in the upper-right quadrant are the anomaly: 2-opt made the
    tour shorter but the finder slower (225 of 1,000 trials; 491 trials
    where 2-opt helped the FCR).}\label{fig:target-delay}
\end{figure}

The experiment supports a named conclusion we call the
\emph{Target-Delay Anomaly}: for finder-centric evaluation
metrics, \emph{greedy probability-weighted chaining dominates
tour-length-optimized chaining}. The greedy chain approximates an
ordering whose prefix is optimized jointly for proximity and prior
mass, which is exactly what the FCR numerator rewards; 2-opt
optimizes the total tour length, a quantity that is a different and
generally orthogonal objective. The result implies that the standard
robotics pipeline of ``build a tour, then polish it with 2-opt'' is
not merely neutral but actively harmful in a multi-robot search
setting with Bayesian priors and finder-centric objectives. It also
provides a concrete reason to prefer the greedy $p/d^2$ chain of
Chapter~III not merely as a real-time simplification but as the
correct ordering under the search objective itself.
